\begin{proposition}[G-11: Explicit rectangular multiplicity gaps]\label{G-11}
For $n\geq4$ define the following positive integers:
\[
\begin{gathered}
 b_n=\binom{n+3}{3},\qquad \beta_n=6b_n n^{30},\qquad
 T_n=\binom{b_n+\beta_n}{b_n}-1,\\
 E_n=4T_n(\beta_n!),\qquad u_n=2n^2-13,\\
 Q_n=(nE_n)^{u_n-1},\qquad K_n=u_n(Q_n-1)+1.
\end{gathered}
\]
For every integer $k\geq K_n$, put $d=E_n k$ and
$\lambda=(nd/4,nd/4,nd/4,nd/4)$. Then
\begin{equation}\label{G11:rectangle}
 m_\lambda(\operatorname{per}_n)>m_\lambda(\det_n).
\end{equation}
More generally, let $n>m\geq4$. The same conclusion, with
$\operatorname{per}_n$ replaced by
$\ell^{n-m}\operatorname{per}_m$, holds if
\begin{equation}\label{G11:padded-hypotheses}
 \dim\mathcal P_{m,4}+3>2n^2+2
 \quad\hbox{and}\quad n<4m/3.
\end{equation}
In particular it holds for $n=m+1$ and every $m\geq8$, and
for $n=\lfloor\rho m\rfloor$, all sufficiently large $m$,
and any fixed $1<\rho<\sqrt5/2$.
The degree, its progression, and the rectangular partition are
explicit. No useful computational size or polynomial algorithm
is asserted for these bounds.
\end{proposition}
\paragraph{Background used.}
We use the Kempf--Ness closed-orbit criterion in Losev (2006),
Theorem 4, p.~1; the stable-locus and quotient statements of
Mukai (2003), Definition 5.12, Proposition 5.15, Theorem 5.16,
and Corollary 5.17, pp.~165--167; and Derksen (2001),
Theorem 1.1, Proposition 1.2, and Example 2.1, printed p.~956.
The latter gives the displayed bound $\beta_n$ for generators
of the invariant ring of quaternary degree-$n$ forms.
The argument below proves the particular stability witnesses,
Veronese bound, and GCT deduction that those sources do not state.
\begin{proof}
Write $G=\mathrm{SL}_4(\mathbb C)$ and let $V_n$ be the space
of degree-$n$ quaternary forms. First we exhibit stable points
in the cones whose dimensions we compare. The form
\[
 f_n=x_0^n+x_1^n+x_2^n+x_3^n
\]
belongs to both $\mathcal P_{n,4}$ and $\mathcal D_{n,4}$.
Indeed factor $x_0^n+x_1^n$ and $x_2^n+x_3^n$ into binary
linear factors, place them on the diagonal and directed cycle
of an $n\times n$ pencil, and use its two perfect matchings.
For the determinant, multiply one cycle entry by the sign
$(-1)^{n-1}$.

Equip $V_n$ with an $\mathrm{SU}_4$-invariant Hermitian inner
product and use unit monomial vectors. Distinct monomials are
orthogonal. In the normalized sum of the four pure powers,
the mean exponent of each variable is $n/4$. Its moment map
therefore vanishes on the diagonal trace-zero Lie algebra.
It also vanishes on the off-diagonal directions: an operator
$x_j\partial/\partial x_i$, $i\ne j$, cannot carry one of
these pure powers to another when $n\geq3$. Thus the moment
map is zero. The Kempf--Ness criterion makes the $G$-orbit
closed. Its infinitesimal stabilizer is zero, since the
monomials $x_i^{n-1}x_j$ and $x_i^n$ in the derivative of
$f_n$ are distinct for $n\geq3$ and force all matrix entries
to vanish. Its stabilizer is consequently finite. This is a
stable point in each cone.

For the padded case put $a=n-m>0$ and suppose $3a<m$.
Let $e_0$ be the unit monomial $x_0^n$ and $e_i$ the unit
monomial proportional to $x_0^a x_i^m$, $1\leq i\leq3$.
Set
\[
 \theta_0=\frac{m-3a}{4m},\qquad
 \theta_1=\theta_2=\theta_3=\frac{n}{4m},\qquad
 F=\sum_{i=0}^3\sqrt{\theta_i}\,e_i.
\]
All four numbers are positive and sum to one. The mean
exponent of $x_0$ is $n\theta_0+3a\theta_1=n/4$, and
that of $x_i$ is $m\theta_i=n/4$. No two supported monomials
differ by a root $e_i-e_j$, because their differences are
multiples of $m\geq4$. The same moment-map calculation
therefore gives zero. We have $F=x_0^a q$, where
$q=\sum_{i=0}^3 c_i x_i^m$ and all $c_i\ne0$.
The preceding binary-factor construction puts $q$ in
$\mathcal P_{m,4}$, so $F\in\mathcal Z_{m,n}$.

Here is a direct finite-stabilizer check. For a linear vector
field $D_A=\sum_i L_i\partial/\partial x_i$, the equation
$D_A(F)=0$, after division by $x_0^{a-1}$, is
\[
 aL_0q+x_0D_Aq=0.
\]
Reduction modulo $x_0$ shows $L_0=d_0x_0$, because $a\ne0$
and $q(0,x_1,x_2,x_3)\ne0$. It follows that
$D_Aq=-ad_0q$. The distinct mixed monomials
$x_i^{m-1}x_j$ force every off-diagonal entry of $A$ to
vanish. The $x_0^m$ coefficient gives $(m+a)d_0=0$,
and the remaining pure powers give every other diagonal
entry zero. Thus the Lie stabilizer is zero and $F$ is
stable by Kempf--Ness.

For every reduced irreducible cone $X$ just considered,
the stable locus is a nonempty open subset and its quotient
fibers are the $15$-dimensional $G$-orbits (Mukai, loc. cit.).
Hence the dimension-of-fibers theorem yields
\begin{equation}\label{G11:quotient-dimension}
 \dim\mathbb C[X]^G=\dim X-15.
\end{equation}
In particular, writing $A_D=\mathbb C[\mathcal D_{n,4}]^G$
and $A_X=\mathbb C[X]^G$, G-06, R-15, and the first
hypothesis in \eqref{G11:padded-hypotheses} imply
\[
 \dim A_D\leq u_n,\qquad \dim A_X\geq u_n+1.
\]
The invariant functor is exact in characteristic zero.
Thus both invariant rings are graded quotients of
$\mathbb C[V_n]^G$, including when $X$ is padded; see also
Mukai, Proposition 5.10 and its proof, pp.~164--165.
Derksen's bound, specialized to $q=4$ and symmetric degree
$n$, says that this ambient invariant ring is generated in
degrees at most $\beta_n=6b_n n^{30}$. The dual convention
for forms makes no difference: the contragredient
transforms the action by an automorphism of $\mathrm{SL}_4$.
We may take at most $T_n$ homogeneous generators, because
the vector space of all coefficient polynomials in $b_n$
variables of positive degree at most $\beta_n$ has dimension
$T_n$.

We check explicitly that the $E_n$th Veronese of each ring
is generated in its first degree. More generally, suppose
an algebra has $r\leq T_n$ homogeneous generators with
positive degrees $d_i\leq\beta_n$, and put
$B_0=\beta_n!$. In a monomial in these generators of total
degree $kE_n$, remove as many pure blocks of
$B_0/d_i$ copies of generator $i$ as possible. Every removed
block has degree $B_0$. The remaining exponents are less
than $B_0/d_i$, so the residual degree is less than
$rB_0\leq T_n B_0$. It is an integer multiple $hB_0$,
because the original and removed degrees are such multiples;
thus $0\leq h<T_n$. There are $4kT_n-h$ pure blocks.
Combine the residual monomial with $4T_n-h$ of these blocks
to form a monomial of degree $E_n=4T_nB_0$, then divide
the other $4(k-1)T_n$ blocks into groups of $4T_n$.
This factors the original monomial into $k$ monomials of
degree $E_n$. Applying this to spanning monomials proves
the claim, also after passing to a quotient.

Consequently $A_D^{(E_n)}$ and $A_X^{(E_n)}$ are standard
graded domains, with dimensions unchanged. Set
$s=\dim\operatorname{Proj}A_D^{(E_n)}\leq u_n-1$.
This projective dimension is positive: block padding embeds
$\mathcal D_{4,4}$, whose dimension is $34$ by G-07,
into $\mathcal D_{n,4}$. By \eqref{G11:quotient-dimension}
the latter invariant ring has dimension at least $19$,
so $s\geq18$.
Choose a basis of its degree-one piece represented by
homogeneous invariant polynomials of degree $E_n$ on $V_n$.
Their composition with the determinant-pencil map is a
collection of homogeneous polynomials of degree $nE_n$.
They give a dominant rational map from
$\mathbb P^{4n^2-1}$ to $\operatorname{Proj}A_D^{(E_n)}$.
Dominance follows from dominance of the pencil map and
surjectivity of the affine quotient (Mukai, Theorem 5.9);
the Veronese morphism has this projective image. G-10's
base-locus-aware image-degree argument gives
\[
 \deg\operatorname{Proj}A_D^{(E_n)}
 \leq(nE_n)^s\leq Q_n.
\]
Chardin's bound in G-10 therefore gives, for every $k\geq1$,
\[
 \dim(A_D)_{E_nk}
 \leq Q_n\binom{k+u_n-1}{u_n-1}.
\]
Linear Noether normalization of the standard graded ring
$A_X^{(E_n)}$, whose dimension is at least $u_n+1$, gives
\[
 \dim(A_X)_{E_nk}\geq\binom{k+u_n}{u_n}.
\]
The ratio of the binomials is $(k+u_n)/u_n>Q_n$ when
$k\geq K_n$. Thus the first invariant dimension is strictly
smaller than the second.

Finally, in coefficient degree $d$, every polynomial
$\mathrm{GL}_4$ weight has size $nd$. The only irreducible
polynomial representation with trivial restriction to
$\mathrm{SL}_4$ is a determinant power; it has shape
$(nd/4)^4$, when $4\mid nd$, and is one-dimensional.
Here $E_n$ is divisible by four. The invariant dimension
difference is therefore exactly the difference of the two
multiplicities for the displayed rectangular partition.
The highest-weight restriction argument of G-05 and R-14
identifies these with the full orbit-closure multiplicities.
This proves \eqref{G11:rectangle} and its padded version.
R-14 and R-16 supply the stated particular ranges; all
satisfy $n<4m/3$.
\end{proof}
\paragraph{Relation to the rectangular bridge.}
The Haar projector onto $\mathrm{SU}_4$ invariants isolates
precisely these rectangular highest weights. Equivalently,
the corresponding character integral is the dimension of
this invariant graded piece. The proof makes that positive
rank exceed the determinant rank; it does not replace an
orbit-closure multiplicity by an ambient Kronecker
coefficient. The quartic formula G-07/R-12 is much sharper
where it applies. These deductions make no literature-priority
claim and are independent of RH and either proposed GRH input.

\begin{proposition}[G-12: A padding exclusion for the rectangular sector]\label{G-12}
Let $m\geq1$, $n>m$, and $n\geq4m/3$. For every degree-$m$
form $f$, every coordinate-ring degree $d\geq1$ with
$4\mid nd$, and $\lambda=(nd/4)^4$, one has
\begin{equation}\label{G12:exclusion}
 m_\lambda(\ell^{n-m}f)\leq m_\lambda(\det_n).
\end{equation}
If $n>4m/3$, the multiplicity on the left is zero.
The equality case $n=4m/3$ is included in the inequality.
Thus the rectangular gaps of G-11 cannot extend to or beyond
this padding boundary, even though nonrectangular four-row
questions remain.
\end{proposition}
\begin{proof}
As in G-05, restrict substituted forms to four variables.
They have shape $L^a q$, where $a=n-m$ and $q$ has degree
$m$. If $L=0$, the form is zero. Otherwise a change of
coordinates takes $L$ to $x_0$; it suffices to use a
$\mathrm{GL}_4$ change, since the determinant cone and each
homogeneous invariant ideal are stable under this group.
Act on the variables by the one-parameter subgroup
\[
 \eta(t)=\operatorname{diag}(t^3,t^{-1},t^{-1},t^{-1})
 \quad\hbox{of }\mathrm{SL}_4.
\]
A monomial in $x_0^a q$ with $x_0$ exponent $a+j$ has
weight $3a-m+4j$. If $3a>m$, all these weights are positive,
so $\eta(t)x_0^a q\longrightarrow0$ as $t\longrightarrow0$.
Every homogeneous positive-degree invariant therefore
vanishes on every such padded form and on their orbit
closure. Exactness of invariants identifies its invariant
coordinate ring with $\mathbb C$, whose positive-degree pieces vanish,
which proves the zero assertion.

If $3a=m$, the same limit instead equals
\[
 F_0=x_0^a q(0,x_1,x_2,x_3).
\]
Every ternary degree-$m$ form belongs to $\mathcal D_{m,3}$
by G-05 (including the elementary degrees one and two).
Embedding it and adding $a$ diagonal copies of $x_0$
shows $F_0\in\mathcal D_{n,4}$; the argument also holds for
limits because the block-padding map is polynomial.
Let $P$ be any homogeneous $\mathrm{SL}_4$ invariant
vanishing on $\mathcal D_{n,4}$. Since $P$ is constant
along $\eta$, continuity gives
$P(x_0^a q)=P(F_0)=0$. Thus the determinant ideal intersected
with the ambient invariant ring is contained in the ideal
of the restricted padded orbit closure intersected with
that ring. Both invariant coordinate rings are quotients
of the same ambient invariant ring, by exactness of
invariants. This containment gives a surjection from the
determinant invariant coordinate ring to the padded one,
and hence the degreewise dimension inequality.
Invariant dimensions are precisely the rectangular
multiplicities, as in G-11; highest-weight restriction
transfers the inequality to the full orbit closures.
\end{proof}
\paragraph{Source and scope.}
The strict inequality is also the rectangular specialization
of the usual padding weight constraint, but the proof above
gives its one-parameter-subgroup form and handles the
boundary by the ternary determinant theorem. This is an
exclusion for this specified rectangular sector. It does
not compare all four-row multiplicities or arbitrary
Kronecker coefficients.
